\section{Single-Pass Compaction Results}
\label{sec:singlepass}

\subsection{Recall degradation is monotone and severe}
\label{sec:degradation}

\begin{figure}[h]
\centering
\includegraphics[width=0.85\linewidth]{fig3_recall_vs_compaction.png}
\caption{Recall degradation by compaction percentage at $Q=5$, for three densities.
All curves decline monotonically. The grep upper bound at $\delta=0.42$ (dashed) stays above 80\%
even at C4, illustrating the grep-LLM divergence.}
\label{fig:recall-compaction}
\end{figure}

Recall degrades monotonically with compaction percentage
(Claude Haiku 4.5 QA, $Q=5$):

\begin{table}[h]
\centering
\caption{Recall vs compaction level across three densities ($Q=5$). Values in parentheses are deltas vs.\ the C0 baseline.}
\label{tab:recall-compaction}
\begin{tabular}{lccc}
\toprule
Level   & $\delta=0.21$ ($Q=5$) & $\delta=0.32$ ($Q=5$) & $\delta=0.42$ ($Q=5$) \\
\midrule
C0       & 62.5\%       & 75.0\%       & 78.8\% \\
C1 (5\%) & 55.0\% ($-7$)  & 68.3\% ($-7$)  & 72.5\% ($-6$) \\
C2 (25\%)& 40.0\% ($-23$) & 50.0\% ($-25$) & 58.8\% ($-20$) \\
C3 (50\%)& 25.0\% ($-38$) & 33.3\% ($-42$) & 43.8\% ($-35$) \\
C4 (98\%)& 2.5\%  ($-60$) & 10.0\% ($-65$) & 7.5\%  ($-71$) \\
\bottomrule
\end{tabular}
\end{table}

Even the lightest compaction (C1, 5\%) costs 6--7 percentage points. At C4
(98\%), recall drops to single digits --- the conversation is effectively
destroyed.

The degradation is consistent across densities: higher-density contexts
have more room to fall but the relative pattern is identical.

\subsection{The compacted zone is dead}
\label{sec:dead-zone}

Facts within the compacted region are almost never recalled, even though
their keywords survive in the summary (Claude Haiku 4.5 QA, $Q=5$, $\delta=0.42$):

\begin{table}[h]
\centering
\caption{Grep keyword survival vs.\ LLM recall in the compacted zone.}
\label{tab:dead-zone}
\begin{tabular}{lcc}
\toprule
Level & Grep survival (compacted zone) & LLM recall (compacted zone) \\
\midrule
C1 & 50\% (n=2)  & 0\% \\
C2 & 82--91\%    & 8--18\% \\
C3 & 83--93\%    & 0--7\% \\
C4 & 59--62\%    & 2--4\% \\
\bottomrule
\end{tabular}
\end{table}

This is the most striking result: \textbf{grep finds the keywords in the summary,
but the LLM cannot use them to answer questions.} The summary preserves
lexical traces of the facts but destroys the surrounding context that would
enable retrieval. This is the ``Lost in the Middle'' effect amplified: the
summary sits at the very beginning of the context --- the least-attended
position after the primacy window is exhausted.

\subsection{Attention dilution: compaction damages untouched facts}
\label{sec:attention-dilution}

The remaining zone (facts that were \textit{not} compacted, sitting in their original
messages) also loses recall as compaction increases:

\begin{table}[h]
\centering
\caption{Remaining-zone recall as compaction increases ($\delta=0.42$, $Q=5$).}
\label{tab:remaining-zone}
\begin{tabular}{lcc}
\toprule
Level & Remaining zone recall ($\delta=0.42$, $Q=5$) & $n_{\text{facts}}$ \\
\midrule
C1 & 73\% & 78 \\
C2 & 59\% & 69 \\
C3 & 53\% & 59 \\
C4 & 80\% & 5 \\
\bottomrule
\end{tabular}
\end{table}

From C1 to C3, remaining-zone recall drops from 73\% to 53\% --- a 20pp loss on
facts that were never touched by compaction. The mechanism: re-padding replaces
the compacted portion with noise (unrelated conversation sessions). This noise
dilutes the model's attention, degrading retrieval even for intact facts.

C4 is an outlier (80\%) because only 5 facts remain in the zone --- too few for
statistical reliability, but consistent with the model finding a needle among
mostly-padding.

\textbf{Caveat --- dilution vs semantic interference.} Our re-padding uses
unrelated \textit{real} conversations rather than information-free filler. The
20pp drop therefore mixes two effects: pure attention dilution (more
tokens to attend to) and semantic interference (the unrelated conversations
may contain content that competes with or partially matches the test
questions). Disentangling them would require a control arm where the
compacted region is replaced with low-information filler (e.g.\ random
tokens or repeated boilerplate). We leave this control to future work
(\S\ref{sec:future}).

\subsection{Spatial recall density}
\label{sec:spatial}

\begin{figure}[h]
\centering
\includegraphics[width=0.85\linewidth]{fig4_spatial_recall.png}
\caption{Cumulative recalled facts by position in the original 190K context ($\delta=0.42$, $Q=5$).
Each step curve shows one compaction level. Dotted vertical lines mark compaction boundaries
(C1 at 9.5K, C2 at 47K, C3 at 95K, C4 at 186K). The C0 baseline rises near-linearly: at the
sweet-spot density, recall is roughly uniform across position. Compacted variants stay close
to flat throughout the compacted zone (low slope) and recover the C0 slope only in the
surviving portion --- yielding the expected two-slope profile, with the slope change located
at each level's compaction boundary.}
\label{fig:spatial-recall}
\end{figure}

This visualization maps each fact to its position in the original (pre-compaction)
context and shows whether it was recalled. It reveals:

\begin{itemize}
  \item \textbf{C0}: Relatively uniform recall across the context, with slight primacy
    and recency effects
  \item \textbf{C1}: A small shadow at the beginning (compacted zone), rest intact
  \item \textbf{C2}: A larger shadow covering the first 25\% of the context
  \item \textbf{C3}: The first half is dead, the second half is weakened by attention dilution
  \item \textbf{C4}: Almost complete darkness, with tiny islands of recall at the very end
\end{itemize}

This provides a visual ``damage map'' showing exactly where compaction destroys
information.

\subsection{Repeatability}
\label{sec:singlepass-repeatability}

Three independent runs at $\delta=0.42$ / $Q=5$ (the most demanding configuration
with 80 facts; Claude Haiku 4.5 QA + Haiku judge) confirm stability:

\begin{table}[h]
\centering
\caption{Three-run repeatability at $\delta=0.42$, $Q=5$.}
\label{tab:repeatability}
\begin{tabular}{lccccc}
\toprule
Level & Run 1 & Run 2 & Run 3 & Mean & $\sigma$ \\
\midrule
C0 & 73.8\% & 78.8\% & 77.5\% & 76.7\% & $\pm 2.6$ \\
C1 & 71.2\% & 68.8\% & 68.8\% & 69.6\% & $\pm 1.4$ \\
C2 & 52.5\% & 53.8\% & 52.5\% & 52.9\% & $\pm 0.7$ \\
C3 & 40.0\% & 40.0\% & 40.0\% & 40.0\% & $\pm 0.0$ \\
C4 & 8.8\%  & 6.2\%  & 6.2\%  & 7.1\%  & $\pm 1.5$ \\
\bottomrule
\end{tabular}
\end{table}

Maximum variance is $\pm 2.6$pp (C0 baseline). All compaction effects ($-7$pp to
$-70$pp) are far larger than measurement noise. The C3 result is remarkably
stable (40.0\% in all three runs), suggesting that at this compaction level,
the outcome is nearly deterministic.

\subsection{Q-effect inversion under compaction}
\label{sec:q-inversion}

The calibration (\S\ref{sec:calibration}) established that asking more questions per prompt yields
higher recall in static contexts: at $\delta=0.42$, $Q=10$ reaches 79\% vs 68\% for $Q=1$.
This holds for uncompacted and lightly compacted contexts. However, under
severe compaction, the effect \textbf{reverses} (Claude Haiku 4.5 QA, $\delta=0.42$):

\begin{table}[h]
\centering
\caption{Q-effect inversion under severe compaction. At C4, $Q=1$ surpasses $Q=5$.}
\label{tab:q-inversion}
\begin{tabular}{lccc}
\toprule
Level & $Q=5$ & $Q=1$ & $\Delta$ ($Q=1 - Q=5$) \\
\midrule
C0       & 73.8\% & 67.5\% & $-6.3$pp \\
C1 (5\%) & 66.2\% & 55.0\% & $-11.2$pp \\
C2 (25\%)& 50.0\% & 47.5\% & $-2.5$pp \\
C3 (50\%)& 35.0\% & 31.2\% & $-3.8$pp \\
C4 (98\%)& 2.5\%  & 11.2\% & \textbf{+8.7pp} \\
\bottomrule
\end{tabular}
\end{table}

At C4, $Q=1$ surpasses $Q=5$ by 8.7pp. The same inversion appears in
the iterative multi-pass experiments (\S\ref{sec:strategy}): across all four strategies at 1M
tokens, $Q=1$ outperforms $Q=5$ by 5--9pp.

The mechanism is intuitive: in a rich context, multiple questions per prompt
cast a wider net, increasing the chance of activating diverse regions. In a
degraded context (heavy compaction or many compaction cycles), the model's
limited attention is better spent focusing on a single question rather than
splitting across multiple retrieval targets.

This has practical implications for benchmarking: \textbf{$Q=5$ results are not
directly comparable to $Q=1$ results}, and the direction of the bias depends
on context quality. Evaluations of compacted contexts should report results
at multiple $Q$ values or, at minimum, note which $Q$ was used.

\subsection{Cross-model validation: Sonnet 4.6}
\label{sec:sonnet}

To test whether our findings are model-specific, we replicated the
single-pass compaction experiment using Claude Sonnet 4.6 as the QA model
($Q=1$, $\delta=0.42$; judge: Haiku 4.5). Both columns below are $Q=1$ --- this is
the configuration in which the cross-model gap is largest, see \S\ref{sec:q-inversion}.

\begin{table}[h]
\centering
\caption{Cross-model comparison: Haiku 4.5 vs Sonnet 4.6 at $Q=1$, $\delta=0.42$.}
\label{tab:sonnet}
\begin{tabular}{lcccc}
\toprule
Level & Haiku 4.5 ($Q=1$) & Sonnet 4.6 ($Q=1$) & Haiku $\Delta$/C0 & Sonnet $\Delta$/C0 \\
\midrule
C0       & 67.5\% & 92.5\% & ---       & --- \\
C1 (5\%) & 55.0\% & 90.0\% & $-12.5$pp & $-2.5$pp \\
C2 (25\%)& 47.5\% & 86.2\% & $-20.0$pp & $-6.2$pp \\
C3 (50\%)& 31.2\% & 62.5\% & $-36.2$pp & $-30.0$pp \\
C4 (98\%)& 11.2\% & 21.2\% & $-56.2$pp & $-71.2$pp \\
\bottomrule
\end{tabular}
\end{table}

Three key findings emerge:

\textbf{1. Sonnet has no Lost-in-the-Middle effect.} The C0 spatial recall profile
is flat: 93\% early, 93\% mid, 92\% late. Haiku's C0 shows the classic U-shaped
profile (higher at extremes, lower in the middle). This difference disappears
under compaction --- both models converge to similar spatial patterns at C3--C4.

\begin{figure}[h]
\centering
\includegraphics[width=0.95\linewidth]{fig8_recall_density.png}
\caption{Smoothed recall density by position in the original context (Gaussian
kernel, bw=15kTok). Left: Haiku $Q=5$. Center: Haiku $Q=1$. Right: Sonnet $Q=1$.
Sonnet's C0 (green) is flat at $\sim$95\%, confirming the absence of Lost-in-the-Middle.
Under compaction, all models show the same pattern: dead zone in the compacted
region, concentration of recall at the context end.}
\label{fig:recall-density-smoothed}
\end{figure}

\textbf{2. A stronger model resists light compaction better.} Sonnet loses only
2.5pp at C1 and 6.2pp at C2, vs 12.5pp and 20.0pp for Haiku. The additional
25pp of baseline recall provides a buffer against moderate information loss.

\textbf{3. Severe compaction equalizes models.} At C3--C4, the delta relative to C0
is similar ($-30$pp and $-71$pp for Sonnet vs $-36$pp and $-56$pp for Haiku). This
confirms that the information is \textbf{destroyed by the compaction process itself},
not merely hidden from a weaker model's attention. If the degradation were
purely attentional, Sonnet's superior baseline should provide proportional
protection at all compaction levels.

This cross-model comparison strengthens the paper's central finding: the
bottleneck at high compaction is information destruction, not attention
capacity.
